\section{模型的假设}
为了简化问题，本文作如下假设：
\begin{enumerate}[label=\textbf{假设\arabic*：},leftmargin=4.5em,labelsep=0.4em,itemsep=0.3em]
    \item 附件 1、2、4 中，每个时刻记录的功率均为此前 10 分钟内的平均功率，电价对应同一时段。
    \item 充电效率和放电效率均为 90\%。充放电功率及电量按交流侧计量，储电量按电池内部计量。
    \item 问题一每日初、末储电量均为 6000 kWh。问题二至四中，1 月储能停用，储电量保持初值；2 月起每日初值等于前一日末值，当日末值随充放电变化。
    \item 各时段功率取该时段平均值，控制系统即时响应，采样、通信和执行延迟均设为0。紧急购电全部用于补足负荷缺口。
    \item 制定计划仅使用决策时已知的数据和已发布的预报。问题四按预测电价制定计划，按实际交付电价结算。
    \item 允许弃用多余的光伏发电量和已购买但未使用的电量。
    \item 问题三在 0:00 确定并支付全天基础购电计划；6:00、12:00 和 18:00 获得新预报后重新规划至当天 24:00，但每次只确认至下一预报发布时刻的购电调整。其后时段用于储能价值评估，交易记录和调整费均对应本次确认窗口。
    \item 问题三和问题四-3按最终确认量相对零点计划结算。增购部分按当期电价的 1.5 倍收费；减购时退回原购电费，并对取消部分支付当期电价 50\% 的违约费。
\end{enumerate}
